\documentclass[output=paper]{langscibook}
\ChapterDOI{10.5281/zenodo.22880264}
\title{\textit{A takinima} – The woven bamboo wall}
\author{Enoch Horai Magum}
\abstract{}
\begin{document}
\maketitle\label{chap:19}
\begin{figure}[H]
    \includegraphics[trim=0 0 0 0,clip,width=.8\linewidth]{figures/19_1.jpg}
    \caption{\textit{A inu gorogorohu} -- A sleeping house}
    \label{fig:19_1}
\end{figure}
\begin{xltabular}{\textwidth}{lXX}
    (\ref{ex:text14-1})&\textit{A takinima eve a peha taba to tabotaboo vai niara nomanai en Teapu.}&Plaited bamboo is a thing that we use today for making walls here in the Teop region.\\
    (\ref{ex:text14-2})&\textit{Eve he na paku raara bona koopu.}&We make it from \textit{koopu} bamboos.\\
    (\ref{ex:text14-3})&\textit{A kiu nae he na paku vonaen:}&The work on it is done as follows:\\
    (\ref{ex:text14-4})&\textit{Eara na tagihu rara rasuu}&We go into the jungle\\
    (\ref{ex:text14-5})&\textit{ei kou a koopu sa teitei vira bata haana kasuana.}&because there are not many \textit{koopu} bamboos along the coast.\\
    (\ref{ex:text14-6})&\textit{A koopu vai a mataa ae bebeahu na teitei nana rasuu.}&The good and long \textit{koopu} bamboos grow in the jungle.\\
    (\ref{ex:text14-7})&\textit{Beara tagihu ei to sun nae a koopu,}&When we go into the jungle where the \textit{koopu} bamboos stand,\\
    (\ref{ex:text14-8})&\textit{eara toro siisira rakaha roho a koopu, eara repaa mosi mana.}&we must intensively examine the \textit{koopu} bamboos, before we eventually cut them.\\
    (\ref{ex:text14-9})&\textit{O koopu toro mosi ara toro manoto bara toro mahaka rakaha mee.}&The \textit{koopu} bamboo that we should cut must be mature and also must be very dry.\\
    (\ref{ex:text14-10})&\textit{Eara sa antee haara tea mosi o koopu vai o kobus,}&We cannot cut the moist \textit{koopu} bamboo,\\
    (\ref{ex:text14-11})&\textit{tabae na antee nana tea koava ge mahamahan,}&because it can rot or leave spaces,\\
    (\ref{ex:text14-12})&\textit{beara takinima nie,}&when we weave the bamboo wall with it,\\
    (\ref{ex:text14-13})&\textit{ei be raka bata maa a koopu.}&that is when the \textit{koopu} bamboos dry up.\\
    (\ref{ex:text14-14})&\textit{Beara sisira vakavara, eara repaa mosi.}&When we have finished examining (them), we cut (them).\\
\end{xltabular}
\begin{figure}[H]
    \captionsetup[subfigure]{oneside,margin={0cm,0cm}}
    \subfloat[]{\includegraphics[trim=0 0 0 0,clip,width=.30\linewidth]{figures/19_2.jpg}\label{fig:19_2}}
    \subfloat[]{\includegraphics[trim=0 0 50 10,clip,width=.365\linewidth]{figures/19_3.jpg}\label{fig:19_3}}
    \subfloat[]{\includegraphics[trim=0 0 0 0,clip,width=.335\linewidth]{figures/19_4.jpg}\label{fig:19_4}}
    \caption{(a) \textit{A dee koopu} -- A load of bamboos, (b) \textit{Kopakopaa koopu} -- Cut bamboos lengthwise, (c) \textit{O komano koopu} -- The inside of the (shaved) bamboos}
\end{figure}
\begin{xltabular}{\textwidth}{lXX}
    (\ref{ex:text14-15})&\textit{Are mosi vaantee a inu, eara repaa paku amaa dee.}&Then we cut as many as are needed for the house, and then we make loads (to carry the \textit{koopu} bamboos home).\\
    (\ref{ex:text14-16})&\textit{A peha dee kou na antee batana bono peho saavun koopu.}&One load consists of ten \textit{koopu} bamboos.\\
    (\ref{ex:text14-17})&\textit{Kavara vai, eara dee voosu vo vaan.}&After that we carry (them) home to the village.\\
    (\ref{ex:text14-18})&\textit{Ei vaan to kahi paku pete raara o hum vaa tea kopakopaa a koopu.}&It is here in the village where we do the next part (of work), the lengthwise cutting of the \textit{koopu} bamboos.\\
    (\ref{ex:text14-19})&\textit{Beara kopakopaa vakavara, eara repaa paku pete o hum vaa tea tasu.}&When we have cut (them) lengthwise, we do the next part of work,  the beating (the bamboo strips).\\
    (\ref{ex:text14-20})&\textit{A tasu ei; beara tasu sirisiri bata rara a koopu, ore tamee teo vuri beara takitakinima batara.}&This beating (means) that we beat and split the \textit{koopu} bamboo strips, so that they are flexible when we are weaving the bamboo wall.\\
    (\ref{ex:text14-21})&\textit{Tea tasu he eara na tasu rara a koopu bona toraara.}&As for the beating, we beat the \textit{koopu} bamboo strips with an axe.\\
    (\ref{ex:text14-22})&\textit{Be kavara a kiu vaa tea tasu, eara repaa taatagi o komano koopu --}&When the work of the beating is finished, we prepare the inside of the \textit{koopu} bamboo strips --\\
    (\ref{ex:text14-23})&\textit{ei eara na taraha vatotobin bata koara amaa nava nae.}&this means, we shave its bars and make the surface even.\\
    (\ref{ex:text14-24})&\textit{Ei to mene ato vakisi bata ni nana bono koopu teo vuri vaa tea takitakinima.}&This is to prevent the bamboo from getting stuck (at the knots) when the weaving is done.\\
    (\ref{ex:text14-25})&\textit{Beara taatagi vakavara a koopu, eara nao vai tea kiu vaa tea takitakinima.}&When we have prepared the \textit{koopu} bamboo strips, we proceed with the work of the bamboo-weaving.\\
    (\ref{ex:text14-26})&\textit{Tea takitakinima he na tei me nana bono buaku o kaku hanana vaa tea paku takinima.}&As for the bamboo-weaving, there are two kinds of patterns of making woven bamboo walls.\\
    (\ref{ex:text14-27})&\textit{O peho o vavaasun, o meho he o vavaagoroho.}&One is the make-stand-up pattern, the other one the make-asleep pattern.\\
\end{xltabular}
\begin{figure}[H]
    \captionsetup[subfigure]{oneside,margin={0cm,0cm}}
    \captionsetup{justification=centering}
    \subfloat[]{\includegraphics[trim=0 0 0 0,clip,width=.5\linewidth]{figures/19_5.jpg}\label{fig:19_5}}\hfil
    \subfloat[]{\includegraphics[trim=0 0 0 0,clip,width=.5\linewidth]{figures/19_6.jpg}\label{fig:19_6}}
    \caption{(a) \textit{O vavaasun} -- The make-stand-up pattern,\\ (b) \textit{O vavaagoroho} -- The make-asleep pattern}
\end{figure}
\begin{figure}[H]
    \captionsetup[subfigure]{oneside,margin={0cm,0cm}}
    \captionsetup{justification=centering}
    \subfloat[]{\includegraphics[trim=0 0 0 0,clip,width=.45\linewidth]{figures/19_7.jpg}\label{fig:19_7}}\hfil
    \subfloat[]{\includegraphics[trim=0 24 0 0,clip,width=.55\linewidth]{figures/19_8.jpg}\label{fig:19_8}}
    \caption{(a) \textit{O vavaasun} -- The make-stand-up pattern,\\ (b) \textit{O vavaagoroho} -- The make-asleep pattern}
\end{figure}
\begin{xltabular}{\textwidth}{lXX}
    (\ref{ex:text14-28})&\textit{Teo vavaasun, eara na ipinoo rara a sun hiava bara a tatabonata na inu.}&As for the make-stand-up pattern, we measure the height and the width of the house.\\
    (\ref{ex:text14-29})&\textit{Eara repaa kurukurus a koopu vai to paa tasu vaevuru ara.}&Then we cut the \textit{koopu} bamboo stripes that we have already beaten (according to the measures).\\
    (\ref{ex:text14-30})&\textit{Be kavara o hum bona, eara repaa vavaasun nao a koopu, ore no suguna komana teo vagegeana na taboo.}&When this part is finished, we make the \textit{koopu} bamboo strips stand up, until the ends of the wall are reached.\\
    (\ref{ex:text14-31})&\textit{Eara repaa taneo tea hogohogo ni nao a koopu vaantee tea tatabonata na taboo.}&Then we start pushing the (horizontal) \textit{koopu} bamboo strips (in between the vertical ones), until it is enough for the width of the wall.\\
    (\ref{ex:text14-32})&\textit{Beara hogo, eara na gae rara a ta kukan, eara repaa hogo tavusu nie.}&When we push them, we lift three (vertical) strips, then we push them (the horizontal ones) through.\\
    (\ref{ex:text14-33})&\textit{Eara repaa taono a ta meha ta kukan, eara repaa gae vaha bana a ta kukan.}&Then we put three other strips on top, and then we lift three strips again.\\
    (\ref{ex:text14-34})&\textit{Eara paku bata koa vo nao raen ore poosu komana a taboo.}&We just do it like this until the wall is completely closed.\\
    (\ref{ex:text14-35})&\textit{Teo vavaagoroho he na vaapeha koana, amaa hogo ae a gae kukan. }&With the make-asleep pattern it is just the same, some pushing and the lifting of three (strips).\\
    (\ref{ex:text14-36})&\textit{Evehee ei kou eara sa paku haara to ipinoo.}&But now we do not do any measurement.\\
    (\ref{ex:text14-37})&\textit{Eara na vaagoroho rara a ta kukan ge a ta taonim,}&We lay three or five strips down,\\
    (\ref{ex:text14-38})&\textit{eara repaa taneo kou tea hogo ae gae, ore suguna komana bea taboo poosu vakavara.}&then we start pushing and lifting (the strips) until the wall is completely closed.\\
    (\ref{ex:text14-39})&\textit{Ahiki to ipinoo, eikou eara na ponuu kaakoo bata koara a koopu beara hogo vaantee bata naora a taboo.}&There is not any measure, because we just bend and move the bamboo strips along, when we are pushing (them) to the (end of the) wall.\\
\end{xltabular}
\begin{figure}[H]
    \includegraphics[trim=0 0 0 0,clip,width=\linewidth]{figures/19_9.jpg}
    \caption{\textit{A inu gorogoroho} -- A sleeping house}
    \label{fig:19_9}
\end{figure}
\end{document}